\section{Substrate model: lattice microphysics and continuum (IR) effective description}
\label{sec:continuum-model}

\subsection{Continuum long-wavelength limit: ontology and kinematics}

Let $\Omega\subset\R^3$ be the intrinsic (material) brane domain.
We use either a periodic cell $\Omega=\T^3$ or a bounded cube with clamped boundary points.
Time $t\in\R$ is an external evolution parameter, not an additional coordinate of the brane.

The fundamental field is the embedding map
\begin{equation}
  \vecX(x,t)\in\R^4,
  \qquad x=(x^1,x^2,x^3)\in\Omega,\quad t\in\R,
  \label{eq:fundamental-field}
\end{equation}
whose components $X^A(x,t)$ ($A=1,\dots,4$) give the brane configuration in the ambient space.

\subsection{Discrete lattice substrate (microphysical postulate)}
\label{subsec:lattice-substrate}

The continuum model is treated as the long-wavelength effective theory of a \emph{cubic lattice} substrate.
Let lattice sites be labeled by $n\in\mathbb{Z}^3$ with spacing $a$ and (reference) material positions
$x_n=a\,n$.
Each site carries an embedded configuration $\vecX_n(t)\in\R^4$.
Interactions are encoded by springs/couplings across neighbor shells with tunable strengths:
nearest neighbors (axis-aligned) with stiffness $k_1$, face diagonals with $k_2$, and body diagonals with $k_3$.
The rest lengths follow the geometry ($a$, $\sqrt{2}\,a$, $\sqrt{3}\,a$) with a uniform prestretch parameter $\alpha$.
The ratios $(k_2/k_1,k_3/k_1)$ can be tuned to suppress leading anisotropic terms in the small-$|k|$ expansion.

\paragraph{Dispersion, anisotropy, and birefringence.}
At wavelengths comparable to $a$, the lattice generically exhibits direction-dependent dispersion and branch
splitting (birefringence), because the cubic symmetry is not fully rotational.
In contrast, for $\lambda\gg a$ (small $|k|$ inside the Brillouin zone), the dispersion can be expanded and the
leading-order dynamics is well-approximated by an isotropic continuum elasticity model with effective Lam\'e-type
parameters.
This motivates using an isotropic hyperelastic action as a baseline in the main text while keeping the lattice
as the microphysical hypothesis.

\paragraph{Brillouin zone and geometric-phase diagnostics.}
The lattice spacing $a$ defines a Brillouin zone $k_i\in[-\pi/a,\pi/a]$.
Because Berry curvature is naturally defined over parameter spaces such as $\mathbf{k}$-space eigenbundles,
the lattice formulation provides a direct route to Berry-curvature diagnostics on a discretized momentum lattice
(see \Cref{sec:geophase}).

\subsection{Induced metric and isotropic reference metric}

The embedding induces a (time-dependent) Riemannian metric on the brane:
\begin{equation}
  g_{ij}(x,t) := \partial_i \vecX(x,t)\cdot \partial_j \vecX(x,t),
  \qquad i,j=1,2,3.
  \label{eq:induced-metric}
\end{equation}
This definition is invariant under rigid motions $\vecX\mapsto Q\vecX+a$ with $Q\in O(4)$ and $a\in\R^4$.
Rotational invariance within the brane is obtained by constructing the stored energy density from invariants of $g$
relative to an isotropic reference metric.

We choose a homogeneous, isotropic reference metric
\begin{equation}
  g^0_{ij}=\ell_0^2\,\delta_{ij},
  \label{eq:reference-metric}
\end{equation}
where $\ell_0>0$ is the continuum analogue of a discrete spring rest length.
Define the mixed relative metric
\begin{equation}
  \widehat g := (g^0)^{-1}g.
  \label{eq:relative-metric}
\end{equation}
Any isotropic hyperelastic energy density can be written as $W=W(\widehat g)$, equivalently as a function of the
principal stretches or invariants of $\widehat g$.

\paragraph{Continuum analogue of pre-tension.}
The reference metric $g^0$ encodes the stress-free geometric scale.
If boundary conditions enforce a held ground-state scale $h_\star$ (uniform), then
\begin{equation}
  \alpha:=\frac{\ell_0}{h_\star}\in(0,1]
  \label{eq:alpha-def}
\end{equation}
measures the mismatch between stress-free scale and held scale.
For $\alpha<1$ the brane is uniformly pre-stressed (continuum ``pre-tension''). This mismatch controls how strongly
transverse gradients backreact on in-brane stresses through the geometric nonlinearity of $g_{ij}(\vecX)$.
In the lattice substrate picture (\Cref{subsec:lattice-substrate}), the same $\alpha$ multiplies the geometric
neighbor distances to set spring rest lengths; \Cref{eq:alpha-def} is the continuum encoding of that uniform prestretch.

\subsection{Constitutive law: St.\ Venant--Kirchhoff (Green--Lagrange strain)}

As a concrete baseline we use the St.\ Venant--Kirchhoff (StVK) energy expressed in terms of the
Green--Lagrange strain tensor
\begin{equation}
  E := \tfrac12\big(\widehat g - I\big)
  \qquad\text{(with $I$ the $3\times 3$ identity)}.
  \label{eq:green-lagrange-strain}
\end{equation}
The StVK stored energy density is
\begin{equation}
  W_{\text{StVK}}(E)
  = \mu\,\mathrm{tr}(E^2)+\frac{\lambda}{2}\,\big(\mathrm{tr}E\big)^2,
  \label{eq:stvk}
\end{equation}
with Lam\'e parameters $\mu>0$ and $\lambda$ (units: Pa). In terms of Poisson ratio $\nu\in(0,1/2)$,
\begin{equation}
\nu=\frac{\lambda}{2(\lambda+\mu)},
\qquad
\lambda=\frac{2\mu\nu}{1-2\nu}.
\label{eq:poisson-lame}
\end{equation}
This choice is not claimed unique; it is used because it is isotropic, analytic, and captures the
geometric nonlinearity through $g_{ij}(\vecX)$.
It also reduces to standard linear elasticity in the small-strain limit, making the longitudinal/transverse
wave speeds and the linearized branch structure transparent.
We emphasize that StVK is employed here as a \emph{baseline constitutive law} for analytic tractability, not as a
claim of material realism at arbitrarily large strains; more robust hyperelastic energies (e.g.\ neo-Hookean,
Mooney--Rivlin, Ogden-type) can be substituted within the same induced-metric framework without changing the
geometric mechanism of transverse--lateral backreaction.

\subsection{Lagrangian, action, and equations of motion}

Let $\rho_m>0$ be the brane mass density (kg/m$^3$). The Lagrangian density is
\begin{equation}
  \Lag(\vecX,\partial_t\vecX,\nabla\vecX)
  =\frac{\rho_m}{2}\,|\partial_t\vecX|^2
  -W\!\left(g(\vecX);g^0\right).
  \label{eq:lagrangian-density}
\end{equation}
The action is
\begin{equation}
  S[\vecX] = \int_{\R}\dd t \int_{\Omega} \dd^3 x\;\Lag.
  \label{eq:action}
\end{equation}

Define the first Piola stress (as a momentum conjugate to $\partial_i X^A$)
\begin{equation}
  P_A^{\;i} := \frac{\partial W}{\partial(\partial_i X^A)}.
  \label{eq:piola-def}
\end{equation}
For energies depending on $g_{ij}$ one can write $P_A^{\;i}=2S^{ij}\partial_j X^A$ with a symmetric
second Piola-type tensor $S^{ij}=\partial W/\partial g_{ij}$.
The Euler--Lagrange equations take divergence form
\begin{equation}
  \rho_m\,\partial_{tt} X^A = \partial_i P_A^{\;i},
  \qquad A=1,\dots,4,
  \label{eq:eom}
\end{equation}
showing explicitly that isotropy (absence of built-in preferred directions) follows from writing
$W$ in terms of invariants of the induced metric relative to the reference metric $g^0$.

\subsection{Spherical material coordinates (coordinate change only)}
\label{subsec:spherical-material-coords}

\paragraph{Important: this is only a coordinate change.}
All dynamical statements remain those of the coordinate-free formulation
$\vecX:\Omega\times\R\to\R^4$ with induced metric $g_{ij}=\partial_i\vecX\cdot\partial_j\vecX$.
We introduce spherical coordinates on the \emph{intrinsic/material} domain as a convenient chart for localized
excitations; this does \emph{not} constrain the embedding $\vecX$.

Let $(r,\theta,\varphi)$ be the standard spherical chart on (a subset of) $\R^3$, with Jacobian
$J(r,\theta,\varphi)=r^2\sin\theta$.
Writing the intrinsic coordinate as $q^a\in\{r,\theta,\varphi\}$ and denoting derivatives in this chart by $\partial_a$,
the equations of motion become
\begin{equation}
  \rho_m\,\partial_{tt}X^A
  =\frac{1}{J(q)}\,\partial_a\!\Big(J(q)\,P_A^{\;a}\Big),
  \qquad A=1,\dots,4,
  \label{eq:eom-spherical}
\end{equation}
which is exactly the Cartesian divergence form \eqref{eq:eom} written with the standard Jacobian factor.

\subsection{Linearization and polarized branches}

Linearize around a homogeneous background embedding $\vecX_\star(x)=(x,0)$ and small displacements
$\vecu(x,t)=\vecX(x,t)-\vecX_\star(x)$.
Under standard regularity assumptions on $W$, one obtains a linear wave system with (in general) multiple
polarization branches, including transverse and longitudinal modes.
For later use, we emphasize two generic features:
(i) different branches may have different characteristic speeds (and may become dispersive at shorter wavelengths),
and (ii) for $\alpha<1$ the pre-stress couples transverse excursions (in the embedding) back into in-brane stresses
through the induced metric \eqref{eq:induced-metric}.
These branches provide the substrate carriers for the geometric-phase and soliton constructions developed below.

\subsection{Transverse--lateral coupling and an emergent ``gravity'' channel}
\label{subsec:gravity-channel}

A distinctive feature of an embedded elastic medium is that transverse gradients change intrinsic distances.
To make this explicit without removing in-brane dynamics, we use a near-identity decomposition around the
homogeneous background embedding $\vecX_\star(x)=(h_\star x,0)$:
\begin{equation}
  \vecX(x,t)=\big(h_\star x^1+\xi^1(x,t),\,h_\star x^2+\xi^2(x,t),\,h_\star x^3+\xi^3(x,t),\,u(x,t)\big),
  \qquad \vec\xi(x,t)\in\R^3.
  \label{eq:near-identity-embedding}
\end{equation}
The induced metric is then
\begin{align}
  g_{ij}
  &= \partial_i\vecX\cdot\partial_j\vecX \nonumber\\
  &= (h_\star\delta_i^{\,a}+\partial_i\xi^a)(h_\star\delta_j^{\,a}+\partial_j\xi^a)+\partial_i u\,\partial_j u \nonumber\\
  &= h_\star^2\delta_{ij}+h_\star(\partial_i\xi_j+\partial_j\xi_i)+\partial_i\vec\xi\cdot\partial_j\vec\xi+\partial_i u\,\partial_j u.
  \label{eq:metric-near-identity}
\end{align}
Setting $\vec\xi\equiv 0$ recovers the familiar graph picture used only for intuition; the ``gravity'' channel,
however, relies on the allowed in-brane relaxation $\vec\xi\neq 0$ sourced by the quadratic term
$\partial_i u\,\partial_j u$.

In the small-slope regime, the in-brane strain has the schematic form
\begin{equation}
  \varepsilon_{ij}\approx \tfrac12(\partial_i\xi_j+\partial_j\xi_i)+\tfrac12\,\partial_i u\,\partial_j u,
  \label{eq:strain-near-identity}
\end{equation}
so the quasi-static balance $\partial_i\sigma_{ij}\approx 0$ with
$\sigma_{ij}=\lambda\,\mathrm{tr}(\varepsilon)\delta_{ij}+2\mu\,\varepsilon_{ij}$ makes it explicit that
the term $\partial_i u\,\partial_j u$ acts as a forcing/eigenstrain-like source for the slow in-brane field
$\vec\xi$.

\paragraph{Pythagorean lengthening and lateral contraction (informal picture).}
For a small material displacement $\dd x$, the ambient line element of the graph embedding satisfies
\begin{equation}
  \dd s^2 = h_\star^2\,|\dd x|^2 + (\nabla u\cdot \dd x)^2.
  \label{eq:line-element-graph}
\end{equation}
Thus any region with nonzero slope $\nabla u$ necessarily \emph{lengthens} intrinsic distances compared to the flat
configuration.
In a discrete spring-lattice realization this is the familiar Pythagorean effect
$|\Delta \vecX|=\sqrt{|\Delta x|^2+(\Delta u)^2}$: transverse variation increases spring length even when lateral
coordinates are unchanged.
In the actual dynamics the in-brane displacement $\vec\xi$ is not constrained; under pre-tension ($\alpha<1$) the
additional metric contribution $\partial_i u\,\partial_j u$ acts as an effective source in the in-brane equilibrium
equations, producing a compensating contraction/dilation pattern rather than a purely local increase of tensile stress.
If the medium is pre-tensioned ($\alpha<1$), the additional length is paid for by increased tensile stress,
which produces net lateral forces that pull material points \emph{toward} the region of transverse excitation.
For a localized bump or rotating transverse pattern, the surrounding brane therefore develops an inward bias
(contraction) rather than an outward pressure.

In this qualitative sense, the ``gravitational mass'' of an excitation is tied to the elastic energy stored in
transverse gradients (and, more generally, to embedded deformation energy), while the ``gravitational field'' is the
slowly varying in-brane strain/dilation pattern sourced by that energy.
The sign is fixed by geometry: increasing transverse arc length increases elastic energy and therefore increases
tension, which pulls laterally inward.

With homogeneous pre-stress ($\alpha<1$), this geometrically induced strain backreacts on the in-brane stress field
through the constitutive law, effectively coupling transverse excitation energy into slow modulations of the
in-brane configuration.

This provides a concrete mechanism for a long-range channel often described informally as ``gravity'' in the present
program: localized transverse gradients can act as sources of slowly varying background strain (contraction/dilation),
and linear wave packets propagating on top of that background experience branch-dependent effective coefficients
(cf.\ \Cref{sec:emergent-relativity}).
In a weak-field, slowly varying regime one may therefore seek a Newtonian-limit reduction in which a scalar potential
built from coarse-grained strain (or elastic energy density) obeys a Poisson-type equation.
Deriving such a reduction quantitatively (and relating the resulting coupling constant to $G$) is an open problem of
the present theory fragment.
